% SOURCE_STATEMENT: "For the time-dependent bivariational initial-value problem used to approximate quantum dynamics, does an approximate trajectory exist, and under what mathematical conditions?"
% Preparation and final self-audit: CONTRIBUTING.md prompts C and D; source/audit artifacts retained in the local physics_statements_2026_10_11 batch.
% !TeX program = lualatex
\documentclass[11pt]{article}
\usepackage[margin=25mm]{geometry}
\usepackage{fontspec}
\setmainfont{Latin Modern Roman}
\usepackage{amsmath,amssymb,mathtools}
\usepackage{unicode-math}
\setmathfont{Latin Modern Math}
\usepackage{xurl,hyperref}
\hypersetup{colorlinks=true,urlcolor=blue}
\setcounter{secnumdepth}{0}
\setlength{\parindent}{0pt}
\setlength{\parskip}{5pt}
\setlength{\emergencystretch}{3em}
\newcommand{\sref}[2]{\hyperlink{source:#1:#2}{[S#2]}}
\newcommand{\cataloguescope}[1]{\paragraph{Recorded target.}#1}
\newcommand{\R}{\mathbb R}
\newcommand{\C}{\mathbb C}
\newcommand{\N}{\mathbb N}
\newcommand{\Z}{\mathbb Z}
\newcommand{\Q}{\mathbb Q}
\newcommand{\F}{\mathbb F}
\newcommand{\eps}{\varepsilon}
\begin{document}
% UnsolvedMath ID 30002892; source catalogue code OWR-13684-002
\setcounter{section}{30002891}
\section{Existence of Time-Dependent Bivariational Quantum Dynamics}\label{30002892}
{\small UnsolvedMath ID 30002892 · Mathematical Physics}
\subsection{Definitions and mathematical statement}

\paragraph{Shared notation.}$\mathbb N=\{1,2,\ldots\}$; the standard number fields have their usual meanings. All additional conventions are stated locally.


\paragraph{Definitions and conventions.}
A bivariational approximation varies a ket $|\psi\rangle$ and an independently parametrized dual bra $\langle\widetilde\psi|$. A typical action has integrand $\langle\widetilde\psi|(i\partial_t-H)|\psi\rangle$; stationarity on the trial manifolds yields coupled evolution equations. The initial-value problem asks for a trajectory through prescribed initial bra and ket data. The trial manifolds, nondegeneracy of their induced pairing and the operator domains control whether this system defines an ordinary differential equation.

\paragraph{Target retained from the canonical source catalogue.}
\paragraph{Catalogue wording retained for traceability.}
For the time-dependent bivariational initial-value problem used to approximate quantum dynamics, does an approximate trajectory exist, and under what mathematical conditions?

\paragraph{Source scope and qualification.}
The action explains the bivariational convention, not an asserted globally well-posed system for arbitrary trial spaces. The question asks existence and conditions, and does not already assume a nondegenerate finite-dimensional manifold.
\subsection{Short English statement}
Under what mathematical hypotheses does time-dependent bivariational quantum approximation produce an actual trajectory?
\subsection{Sources}
\begin{itemize}
\item[\hypertarget{source:30002892:1}{S1}] ULAM.AI and the UnsolvedMath Contributors, UnsolvedMath, record 30002892 (OWR-13684-002). Catalogue curation under CC BY 4.0; definitions and English exposition added locally.
\url{https://huggingface.co/datasets/ulamai/UnsolvedMath}.
\item[\hypertarget{source:30002892:2}{S2}] Mathematical Methods in Quantum Molecular Dynamics (2015), Oberwolfach Reports, DOI 10.4171/owr/2015/27, pp. 1499--1569. Primary report accessed 2026-10-11; relevant PDF pages 61, 62, 63.
\url{https://doi.org/10.4171/owr/2015/27}.
\end{itemize}
\label{end:30002892}
\end{document}
